\begin{longtable}{|c|p{3cm}|p{4cm}|p{4.5cm}|}
\caption{Analisis Risiko dan Strategi Mitigasi} \label{tab:risiko} \\
\hline
\textbf{No} & \textbf{Risiko} & \textbf{Dampak} & \textbf{Strategi Mitigasi} \\
\hline
\endfirsthead
\multicolumn{4}{c}%
{\tablename\ \thetable\ -- \textit{Lanjutan dari halaman sebelumnya}} \\
\hline
\textbf{No} & \textbf{Risiko} & \textbf{Dampak} & \textbf{Strategi Mitigasi} \\
\hline
\endhead
\hline \multicolumn{4}{r}{\textit{Bersambung ke halaman berikutnya}} \\
\endfoot
\hline
\endlastfoot

1 & Keterbatasan data historis untuk pelatihan model peramalan & Akurasi model peramalan tidak mencapai target yang ditetapkan & Menggunakan teknik augmentasi data, transfer learning dari model pre-trained, atau menyesuaikan horizon prediksi sesuai ketersediaan data \\
\hline

2 & Ketidakseimbangan distribusi kelas \textit{intent} & Performa klasifikasi buruk pada kelas minoritas & Menerapkan teknik \textit{oversampling} (SMOTE), \textit{undersampling}, atau \textit{class weighting} pada proses pelatihan \\
\hline

3 & Variasi bahasa alami pengguna di luar cakupan pola yang didefinisikan & Tingkat \textit{fallback} tinggi dan pengalaman pengguna menurun & Memperluas cakupan pola secara iteratif berdasarkan log percakapan dan umpan balik pengguna \\
\hline

4 & Keterbatasan sumber daya komputasi untuk pelatihan model & Waktu pelatihan model melebihi jadwal yang direncanakan & Menggunakan layanan \textit{cloud computing} dengan GPU, mengoptimalkan arsitektur model, atau menggunakan teknik \textit{transfer learning} \\
\hline

5 & Kendala akses ke \textit{Enterprise Data Warehouse} perusahaan & Data pengujian tidak representatif dengan kondisi produksi & Menggunakan data sintetis atau data anonim yang mewakili karakteristik data nyata untuk pengembangan dan pengujian \\
\hline

6 & Perubahan kebutuhan pengguna selama pengembangan & Fitur yang dikembangkan tidak sesuai ekspektasi pengguna akhir & Menerapkan pendekatan iteratif dengan evaluasi berkala bersama pemangku kepentingan dan penyesuaian prioritas \\
\hline

7 & Ketergantungan pada pustaka atau \textit{framework} pihak ketiga & Ketidakcocokan versi atau fitur yang sudah tidak didukung menghambat pengembangan & Mendokumentasikan versi dependensi, menggunakan \textit{virtual environment}, dan menyiapkan alternatif pustaka \\
\hline

8 & Keterbatasan akses responden validasi internal & Evaluasi kebergunaan hanya dapat dilakukan oleh kelompok pengguna terbatas & Mengutamakan validasi mendalam dari tim \textit{Data Management} sebagai pengguna kunci, dan memperluas cakupan validasi secara bertahap \\
\hline

\end{longtable}
